\begin{proof}[Proof of \cref{obj:lem:jackson-packet-localization}]
\begin{enumerate}
\item We first establish the coefficient representation needed for the center estimate. For integers \(N\geq1\) and real \(u\), write
\[
\mathcal R_N(u)=
\begin{cases}
\left\{\dfrac{\sin(Nu/2)}{\sin(u/2)}\right\}^{4},
& \sin(u/2)\neq0,\\[1ex]
N^4,&\sin(u/2)=0.
\end{cases}
\]
The sine addition formulas give
\[
\sin^2((N+1)u/2)-\sin^2(Nu/2)
=\sin((2N+1)u/2)\sin(u/2),
\]
and, for every positive integer \(j\),
\[
\sin((2j+1)u/2)-\sin((2j-1)u/2)
=2\cos(ju)\sin(u/2).
\]
Summing the latter identity, then using the former, yields
\[
\sqrt{\mathcal R_{N+1}(u)}-\sqrt{\mathcal R_N(u)}
=1+2\sum_{j=1}^{N}\cos(ju)
\qquad\text{when }\sin(u/2)\neq0.
\]
When \(\sin(u/2)=0\), one has \(\cos(u/2)=1\) or \(-1\), and the same identity reads
\((N+1)^2-N^2=1+2N\). Starting from \(\mathcal R_1(u)=1\), induction therefore gives
\[
\sqrt{\mathcal R_N(u)}
=N+2\sum_{j=1}^{N-1}(N-j)\cos(ju).
\]

For \(N\geq2\) and \(j\in\mathbb Z\), define
\[
\tau_{N,j}=\max\{N-|j|,0\},\qquad
\gamma_{N,j}=\sum_{j_1\in\mathbb Z}
\tau_{N,j_1}\tau_{N,j-j_1},\qquad
\mathcal S_N=\{j\in\mathbb Z:|j|\leq2N-2\}.
\]
These sums are finite. The sequence \(\tau_{N,j}\) is even and vanishes for
\(|j|\geq N\), so \(\gamma_{N,j}\) is even, nonnegative, and vanishes outside
\(\mathcal S_N\). The preceding expansion becomes
\[
\sqrt{\mathcal R_N(u)}
=\sum_{j\in\mathbb Z}\tau_{N,j}\cos(ju).
\]
To square it, use
\[
\cos(j_1u)\cos(j_2u)
=\frac{\cos((j_1+j_2)u)+\cos((j_1-j_2)u)}2.
\]
Evenness makes the two resulting double sums equal after replacing
\(j_2\) by \(-j_2\). Reindexing their common value by \(j=j_1+j_2\) gives
\[
\mathcal R_N(u)
=\sum_{j_1,j_2\in\mathbb Z}
\tau_{N,j_1}\tau_{N,j_2}\cos((j_1+j_2)u)
=\sum_{j\in\mathcal S_N}\gamma_{N,j}\cos(ju).
\]
At frequency zero,
\[
\gamma_{N,0}
=N^2+2\sum_{j=1}^{N-1}(N-j)^2
=\frac{2N^3+N}{3}>0,
\]
where the last equality follows from the finite sum of squares. Integrating the cosine expansion consequently gives
\[
\int_{-\pi}^{\pi}\mathcal R_N(u)\,\frac{du}{2\pi}
=\gamma_{N,0},
\qquad
c_N=\gamma_{N,0}^{-1}.
\]
Define, for every \(j\in\mathbb Z\),
\[
f_{N,j}=\frac{\gamma_{N,j}}{\gamma_{N,0}}.
\]
Thus
\[
\mathcal J_N(u)=\sum_{j\in\mathcal S_N}f_{N,j}\cos(ju),
\qquad
f_{N,j}\geq0,\qquad
f_{N,-j}=f_{N,j},\qquad
f_{N,j}=0\quad(j\notin\mathcal S_N).
\]

For completeness, cosine orthogonality reads, for \(j_1,j\in\mathbb Z\),
\[
\int_{-\pi}^{\pi}\cos(j_1u)\cos(ju)\,\frac{du}{2\pi}
=
\begin{cases}
1,&j_1=j=0,\\
\frac12,&|j_1|=|j|\neq0,\\
0,&\text{otherwise}.
\end{cases}
\]
This follows by the cosine product identity and integration of the integer modes. The sine part of the Fourier integral vanishes because \(\mathcal J_N\) is even. Hence
\[
\widehat{\mathcal J_N}(j)
=
\begin{cases}
f_{N,0},&j=0,\\
\dfrac{f_{N,j}+f_{N,-j}}2,&j\neq0,
\end{cases}
=f_{N,j}.
\]
In particular, the zero-extended sequences in the statement satisfy
\[
a_j^\pm=
\begin{cases}
f_{N,j}(L\pm j)^{-2},&j\in\mathcal S_N,\\
0,&j\notin\mathcal S_N.
\end{cases}
\]
% lean: CausalSmith.Stat.OptvalueVanishingoverlapRate.kernelCoeff_eq_staged

\item We now obtain the center bounds. Summing the triangular sequence gives
\[
\sum_{j\in\mathbb Z}\tau_{N,j}
=N+2\sum_{j=1}^{N-1}(N-j)=N^2.
\]
Reindexing the finite convolution sum yields
\[
\sum_{j\in\mathbb Z}\gamma_{N,j}=N^4,
\qquad
\sum_{j\in\mathbb Z}|\gamma_{N,j}|
\leq
\sum_{j,j_1\in\mathbb Z}
\tau_{N,j_1}\tau_{N,j-j_1}
=N^4.
\]
Using the value of \(\gamma_{N,0}\), we obtain
\[
\frac N2
\leq
\sum_{j\in\mathcal S_N}f_{N,j}
=\frac{N^4}{\gamma_{N,0}}
\leq2N.
\]
Indeed, the lower inequality is equivalent to \(4N^2\geq1\), which holds for \(N\geq2\), while
\(N^4/\gamma_{N,0}\leq2N\) follows directly from
\(\gamma_{N,0}=(2N^3+N)/3\).

For \(j\in\mathcal S_N\), both shifted frequencies obey
\[
2N\leq L+j\leq6N,\qquad
2N\leq L-j\leq6N.
\]
Their reciprocal squares therefore satisfy
\[
\frac1{36N^2}\leq\frac1{(L\pm j)^2}\leq\frac1{4N^2},
\]
and nonnegativity of \(f_{N,j}\) gives
\[
\frac{f_{N,j}}{18N^2}
\leq a_j^++a_j^-
\leq\frac{f_{N,j}}{2N^2}.
\]
At \(u=0\), every cosine equals one and every coefficient is nonnegative, so
\[
|G_N(0)|=\frac12\sum_{j\in\mathcal S_N}(a_j^++a_j^-).
\]
Consequently,
\[
\frac1{72N}
=
\frac12\,\frac{N/2}{18N^2}
\leq |G_N(0)|
\leq
\frac12\,\frac{2N}{2N^2}
=\frac1{2N}.
\]
% lean: Causalean.Mathlib.Analysis.Approximation.Chebyshev.Jackson.packet_center_bounds

\item To prepare the pointwise estimate, we bound the full-line second differences of the weighted coefficients. Directly from the triangular sequence,
\[
(\delta\tau_N)_j=
\begin{cases}
1,&-N\leq j\leq-1,\\
-1,&0\leq j\leq N-1,\\
0,&\text{otherwise},
\end{cases}
\qquad
\sum_{j\in\mathbb Z}|(\delta\tau_N)_j|=2N.
\]
Shifting the finite convolution sums gives
\[
(\delta\gamma_N)_j
=\sum_{j_1\in\mathbb Z}
(\delta\tau_N)_{j_1}\tau_{N,j-j_1},
\qquad
(\delta^2\gamma_N)_j
=\sum_{j_1\in\mathbb Z}
(\delta\tau_N)_{j_1}(\delta\tau_N)_{j-j_1}.
\]
The triangle inequality, followed by reindexing \(j-j_1\), now yields
\[
\begin{aligned}
\sum_j|(\delta\gamma_N)_j|
&\leq
\sum_{j,j_1}|(\delta\tau_N)_{j_1}|\tau_{N,j-j_1}
=2N^3,\\
\sum_j|(\delta^2\gamma_N)_j|
&\leq
\sum_{j,j_1}|(\delta\tau_N)_{j_1}|
                 |(\delta\tau_N)_{j-j_1}|
=4N^2.
\end{aligned}
\]
Division by the positive constant \(\gamma_{N,0}\) commutes with differences. Thus
\[
\sum_j|f_{N,j}|\leq2N,\qquad
\sum_j|(\delta f_N)_j|
\leq\frac{2N^3}{\gamma_{N,0}}\leq4,\qquad
\sum_j|(\delta^2f_N)_j|
\leq\frac{4N^2}{\gamma_{N,0}}\leq\frac6N.
\]

For \(\varsigma\in\{-1,1\}\) and \(j\in\mathbb Z\), define
\[
\omega_j^{\varsigma}=L+\varsigma j,\qquad
\rho_j^{\varsigma}=
\begin{cases}
(\omega_j^{\varsigma})^{-2},&\omega_j^{\varsigma}\neq0,\\
0,&\omega_j^{\varsigma}=0,
\end{cases}
\qquad
a_j^{(\varsigma)}=
\begin{cases}
a_j^+,&\varsigma=1,\\
a_j^-,&\varsigma=-1.
\end{cases}
\]
The convention at a zero denominator makes \(\rho^{\varsigma}\) a sequence on all of \(\mathbb Z\). Since \(f_N\) vanishes outside \(\mathcal S_N\) and the denominators are positive on \(\mathcal S_N\),
\[
a_j^{(\varsigma)}=f_{N,j}\rho_j^{\varsigma}
\qquad(j\in\mathbb Z).
\]
If \(|j|\leq2N\), then \(N\geq2\) ensures
\[
\omega_j^{\varsigma}\geq N,\qquad
\omega_{j+1}^{\varsigma}\geq N,\qquad
\omega_{j+2}^{\varsigma}\geq N,
\qquad
\omega_{j+1}^{\varsigma}-\omega_j^{\varsigma}=\varsigma.
\]
Factoring the first difference gives
\[
\begin{aligned}
|(\delta\rho^{\varsigma})_j|
&=
\frac1{\omega_j^{\varsigma}(\omega_{j+1}^{\varsigma})^2}
+\frac1{(\omega_j^{\varsigma})^2\omega_{j+1}^{\varsigma}}\\
&\leq\frac2{N^3}\leq\frac{16}{N^3}.
\end{aligned}
\]
For the second difference, the unit spacing of the three denominators gives the exact identity
\[
(\delta^2\rho^{\varsigma})_j
=
\frac{2\{3(\omega_{j+1}^{\varsigma})^2-1\}}
 {(\omega_j^{\varsigma})^2
  (\omega_{j+1}^{\varsigma})^2
  (\omega_{j+2}^{\varsigma})^2}.
\]
Its numerator is nonnegative, and hence
\[
|(\delta^2\rho^{\varsigma})_j|
\leq
\frac6{(\omega_j^{\varsigma})^2(\omega_{j+2}^{\varsigma})^2}
\leq\frac6{N^4}\leq\frac{64}{N^4},
\qquad
|\rho_{j+1}^{\varsigma}|\leq\frac1{N^2}.
\]

The forward-difference product identity is
\[
\begin{aligned}
(\delta^2a^{(\varsigma)})_j
={}&f_{N,j+2}(\delta^2\rho^{\varsigma})_j\\
&+\{(\delta f_N)_{j+1}+(\delta f_N)_j\}
                  (\delta\rho^{\varsigma})_j
+\rho_{j+1}^{\varsigma}(\delta^2f_N)_j.
\end{aligned}
\]
For \(|j|\leq2N\), the preceding estimates imply
\[
\begin{aligned}
|(\delta^2a^{(\varsigma)})_j|
\leq{}&
\frac{64}{N^4}|f_{N,j+2}|
+\frac{16}{N^3}|(\delta f_N)_{j+1}|\\
&+\frac{16}{N^3}|(\delta f_N)_j|
+\frac1{N^2}|(\delta^2f_N)_j|.
\end{aligned}
\]
For \(|j|>2N\), all three values \(f_{N,j}\), \(f_{N,j+1}\), and
\(f_{N,j+2}\) vanish, so this inequality also holds there. In particular,
\(\delta^2a^{(\varsigma)}\) has finite support. Summing and using invariance of full-line sums under integer shifts gives
\[
\begin{aligned}
\sum_j|(\delta^2a^{(\varsigma)})_j|
&\leq
\frac{64}{N^4}\sum_j|f_{N,j}|
+\frac{32}{N^3}\sum_j|(\delta f_N)_j|
+\frac1{N^2}\sum_j|(\delta^2f_N)_j|\\
&\leq\frac{262}{N^3}
\leq\frac{512}{N^3}.
\end{aligned}
\]
% lean: Causalean.Mathlib.Analysis.Approximation.Chebyshev.Jackson.weightedCoeff_delta2_l1

\item We next derive the pointwise localization bound. The reciprocal-square bound on the coefficient support gives, for either sign,
\[
\sum_{j\in\mathbb Z}|a_j^{(\varsigma)}|
\leq\frac1{N^2}\sum_j|f_{N,j}|
\leq\frac2N.
\]
For \(\varsigma\in\{-1,1\}\) and \(u\in\mathbb R\), define the finite Fourier sums
\[
\mathcal F_N^{\varsigma}(u)
=\sum_{j\in\mathbb Z}a_j^{(\varsigma)}e^{iju}.
\]
The definition of \(G_N\) gives
\[
G_N(u)
=-\frac12\operatorname{Re}
\left[
e^{iLu}\mathcal F_N^1(u)
+e^{iLu}\mathcal F_N^{-1}(-u)
\right].
\]
Since the modulation has modulus one,
\[
|G_N(u)|
\leq
\frac{|\mathcal F_N^1(u)|+|\mathcal F_N^{-1}(-u)|}{2}
\leq\frac2N.
\]

For real \(u\), introduce
\[
\eta(u)=\|u\|_{\mathbb T},\qquad
\jmath(u)=\left\lfloor\frac{u}{2\pi}+\frac12\right\rfloor\in\mathbb Z,
\qquad
\upsilon(u)=u-2\pi\jmath(u).
\]
The defining property of the nearest integer gives
\[
|\upsilon(u)|\leq\pi,
\qquad
0\leq\eta(u)\leq|\upsilon(u)|\leq\pi.
\]
Also, replacing the integer in the defining infimum by its negative shows
\[
\eta(-u)=\eta(u).
\]

If \(N\eta(u)\leq1\), then
\[
(1+N\eta(u))^2\leq4,
\qquad
|G_N(u)|\leq\frac2N
\leq\frac{8192}{N(1+N\eta(u))^2}.
\]
This case includes \(\eta(u)=0\).

Suppose instead that \(N\eta(u)>1\), so \(\eta(u)>0\). For every integer
\(j_0\), shifting the finite Fourier sum gives
\[
\sum_j a_{j+j_0}^{(\varsigma)}e^{iju}
=e^{-ij_0u}\mathcal F_N^{\varsigma}(u).
\]
Applying this with shifts one and two yields
\[
\sum_j(\delta^2a^{(\varsigma)})_je^{iju}
=(e^{-iu}-1)^2\mathcal F_N^{\varsigma}(u).
\]
Therefore,
\[
|e^{-iu}-1|^2\,|\mathcal F_N^{\varsigma}(u)|
\leq\sum_j|(\delta^2a^{(\varsigma)})_j|
\leq\frac{512}{N^3}.
\]
Periodicity and the elementary sine bound on \([-\pi/2,\pi/2]\) give
\[
|e^{-iu}-1|
=2|\sin(\upsilon(u)/2)|
\geq\frac2\pi|\upsilon(u)|
\geq\frac2\pi\eta(u)
\geq\frac{\eta(u)}2,
\]
where the last inequality uses \(\pi\leq4\). It follows that
\[
|\mathcal F_N^{\varsigma}(u)|
\leq\frac{2048}{N^3\eta(u)^2}.
\]
Applying this estimate at \(u\) and \(-u\) in the representation of \(G_N\) gives
\[
|G_N(u)|\leq\frac{2048}{N^3\eta(u)^2}.
\]
Finally, \(N\eta(u)>1\) implies
\[
(1+N\eta(u))^2\leq4(N\eta(u))^2,
\]
and hence
\[
|G_N(u)|
\leq\frac{2048}{N^3\eta(u)^2}
\leq\frac{8192}{N(1+N\eta(u))^2}.
\]
Together, the two cases establish
\[
|G_N(u)|
\leq\frac{8192}{N(1+N\|u\|_{\mathbb T})^2}
\qquad(u\in\mathbb R).
\]
% lean: Causalean.Mathlib.Analysis.Approximation.Chebyshev.Jackson.packet_pointwise_localization

\item We obtain the normalized integral bound by integrating a rational envelope. First, for \(|u|\leq\pi\), the integer zero in the defining infimum gives
\(\eta(u)\leq|u|\). For every \(j\leq-1\),
\[
|u-2\pi j|\geq u+2\pi\geq\pi\geq|u|,
\]
and for every \(j\geq1\),
\[
|u-2\pi j|\geq2\pi-u\geq\pi\geq|u|.
\]
The remaining case \(j=0\) gives equality. Thus
\[
\eta(u)=|u|\qquad(-\pi\leq u\leq\pi).
\]
Since
\[
1+(Nu)^2\leq(1+N|u|)^2,
\]
the pointwise estimate implies
\[
\frac{|G_N(u)|}{2\pi}
\leq
\frac{8192}{2\pi N}\,
\frac1{1+(Nu)^2}
\qquad(-\pi\leq u\leq\pi).
\]
The rational envelope has integral
\[
\begin{aligned}
\int_{-\pi}^{\pi}\frac{du}{1+(Nu)^2}
&=\frac{\arctan(N\pi)-\arctan(-N\pi)}N\\
&\leq\frac\pi N,
\end{aligned}
\]
because the arctangent takes values between \(-\pi/2\) and \(\pi/2\).
Both integrands in the comparison are continuous on the compact interval.
Consequently,
\[
\int_{-\pi}^{\pi}|G_N(u)|\,\frac{du}{2\pi}
\leq
\frac{8192}{2\pi N}\,\frac\pi N
=\frac{4096}{N^2}.
\]
% lean: Causalean.Mathlib.Analysis.Approximation.Chebyshev.Jackson.packet_l1_localization

\item All these estimates hold uniformly for \(N\geq2\). Choose
\[
c=\frac1{72},\qquad
C_{\mathrm{center}}=\frac12,\qquad
C_{\mathrm{point}}=8192,\qquad
C_{\mathrm{integral}}=4096,
\]
and set
\[
C=\max\left\{
512,\,
\max\left\{
C_{\mathrm{center}},
\max\{C_{\mathrm{point}},C_{\mathrm{integral}}\}
\right\}
\right\}
=8192.
\]
Then
\[
0<c\leq C,\qquad
512\leq C,\qquad
C_{\mathrm{center}}\leq C,\qquad
C_{\mathrm{point}}\leq C,\qquad
C_{\mathrm{integral}}\leq C.
\]
We now fix any integer \(N\geq2\) and verify the assertions with these common constants.
% lean: CausalSmith.Stat.OptvalueVanishingoverlapRate.jackson_packet_localization

\item The positive normalizing denominator computed above gives
\[
\int_{-\pi}^{\pi}\mathcal J_N(u)\,\frac{du}{2\pi}
=c_N\int_{-\pi}^{\pi}\mathcal R_N(u)\,\frac{du}{2\pi}
=\gamma_{N,0}^{-1}\gamma_{N,0}
=1.
\]
% lean: CausalSmith.Stat.OptvalueVanishingoverlapRate.packetKernel_normalized_integral

\item For either sign, the zero extension implies
\[
(\delta^2a^\pm)_j=0\qquad(|j|>2N).
\]
Hence the absolute second differences are summable, and their full-line sums satisfy
\[
\sum_{j\in\mathbb Z}|(\delta^2a^\pm)_j|
=
\sum_{j=-2N}^{2N}|(\delta^2a^\pm)_j|
\leq\frac{512}{N^3}
\leq\frac C{N^3}.
\]
Because the differences are taken after extending the coefficients by zero, these sums include the differences crossing both endpoints of \(\mathcal S_N\).
% lean: CausalSmith.Stat.OptvalueVanishingoverlapRate.packetCoeffPlus_secondDiff_summable
% lean: CausalSmith.Stat.OptvalueVanishingoverlapRate.packetCoeffMinus_secondDiff_summable
% lean: Causalean.Mathlib.Analysis.Approximation.Chebyshev.Jackson.weightedCoeff_delta2_l1

\item Multiplying the cosine expansion of \(\mathcal J_N\) by \(\cos(Lu)\) gives
\[
k_N(u)
=\frac12\sum_{j\in\mathcal S_N}f_{N,j}
\left\{\cos((L+j)u)+\cos((L-j)u)\right\}.
\]
Every frequency appearing here satisfies
\[
2N+2\leq L\pm j\leq6N-2.
\]
Let \(r\in\mathbb Z\) satisfy the stated alternative. If \(|r|<2N+2\), both shifted frequencies exceed \(|r|\); if \(|r|>6N-2\), both are smaller than \(|r|\). Thus
\[
|L+j|\neq|r|,\qquad |L-j|\neq|r|
\qquad(j\in\mathcal S_N).
\]
Cosine orthogonality consequently gives
\[
\int_{-\pi}^{\pi}\cos((L\pm j)u)\cos(ru)\,du=0.
\]
For the sine integrals, the product identity gives
\[
\begin{aligned}
&\int_{-\pi}^{\pi}\cos((L\pm j)u)\sin(ru)\,du\\
&\qquad=
\frac12\int_{-\pi}^{\pi}
\left\{\sin((r+L\pm j)u)+\sin((r-(L\pm j))u)\right\}\,du
=0.
\end{aligned}
\]
Here every integer sine mode has integral zero; the zero mode is identically zero, and the other modes have equal cosine values at the two endpoints in their antiderivatives. Linearity in the finite expansion therefore yields
\[
\int_{-\pi}^{\pi}k_N(u)\cos(ru)\,du=0,
\qquad
\int_{-\pi}^{\pi}k_N(u)\sin(ru)\,du=0.
\]
% lean: Causalean.Mathlib.Analysis.Approximation.Chebyshev.Jackson.packet_shifted_support

\item Termwise differentiation of the finite cosine polynomial gives, for every real \(u\),
\[
G_N'(u)
=\frac12\sum_{j\in\mathcal S_N}
\left\{
a_j^+(L+j)\sin((L+j)u)
+a_j^-(L-j)\sin((L-j)u)
\right\},
\]
and then
\[
G_N''(u)
=\frac12\sum_{j\in\mathcal S_N}
\left\{
a_j^+(L+j)^2\cos((L+j)u)
+a_j^-(L-j)^2\cos((L-j)u)
\right\}.
\]
The shifted frequencies are positive on \(\mathcal S_N\), so their squares cancel the denominators:
\[
a_j^+(L+j)^2=f_{N,j},
\qquad
a_j^-(L-j)^2=f_{N,j}.
\]
Using
\[
\cos((L+j)u)+\cos((L-j)u)
=2\cos(ju)\cos(Lu),
\]
we conclude that
\[
G_N''(u)
=
\left(\sum_{j\in\mathcal S_N}f_{N,j}\cos(ju)\right)\cos(Lu)
=\mathcal J_N(u)\cos(Lu)
=k_N(u).
\]
% lean: Causalean.Mathlib.Analysis.Approximation.Chebyshev.Jackson.packet_twice_deriv

\item Every frequency \(L\pm j\), for \(j\in\mathcal S_N\), is a nonzero integer. Its cosine integral vanishes by direct integration. Thus
\[
\begin{aligned}
\int_{-\pi}^{\pi}G_N(u)\,du
=-\frac12\sum_{j\in\mathcal S_N}
\left\{
a_j^+\int_{-\pi}^{\pi}\cos((L+j)u)\,du
+a_j^-\int_{-\pi}^{\pi}\cos((L-j)u)\,du
\right\}
=0.
\end{aligned}
\]
% lean: Causalean.Mathlib.Analysis.Approximation.Chebyshev.Jackson.packet_zero_mean

\item Finally, enlarging the three upper constants to the common constant \(C\) gives
\[
\frac cN\leq|G_N(0)|
\leq\frac{C_{\mathrm{center}}}N
\leq\frac CN,
\]
\[
|G_N(u)|
\leq
\frac{C_{\mathrm{point}}}
 {N(1+N\|u\|_{\mathbb T})^2}
\leq
\frac C{N(1+N\|u\|_{\mathbb T})^2}
\qquad(u\in\mathbb R),
\]
and
\[
\int_{-\pi}^{\pi}|G_N(u)|\,\frac{du}{2\pi}
\leq\frac{C_{\mathrm{integral}}}{N^2}
\leq\frac C{N^2}.
\]
These bounds, together with the preceding identities and summability statements, establish every assertion.
% lean: CausalSmith.Stat.OptvalueVanishingoverlapRate.jackson_packet_localization
\end{enumerate}
\end{proof}
